\section{Classical local theory and continuation}\label{sec:local-appendix}
We derive Proposition~\ref{prop:local} from the classical forced local
theory in Tao~\cite[Theorems~5.1 and~5.4]{tao2013}. The reductions below
specify the viscosity, spatial mean, pressure, and continuation conventions
needed here. These results are independent of the compact blowup input.

\begin{lemma}[Sobolev multiplication and embeddings]\label{lem:calculus}
For each integer $m\ge2$ and either domain,
\begin{equation}\label{eq:Rproduct}
 \norm{vw}_{H^m}\le C_m\bigl(\norm{v}_{H^2}\norm{w}_{H^m}
                         +\norm{w}_{H^2}\norm{v}_{H^m}\bigr),
 \qquad \norm{v}_\infty\le C\norm{v}_{H^2}.
\end{equation}
In particular, for integers $k\ge3$,
\begin{align}
 \norm{vw}_{H^k}&\le C_k\norm{v}_{H^k}\norm{w}_{H^k},\label{eq:algebra}\\
 \norm{u\otimes u}_{H^k}&\le C_k\norm{u}_{H^2}\norm{u}_{H^k}.
 \label{eq:tame}
\end{align}
Also $H^1\hookrightarrow L^6$. The homogeneous estimates used in the
critical energy arguments are
\begin{equation}\label{eq:embeddings}
 \norm{v}_3\le C\norm{v}_{\dot H^{1/2}},\qquad
 \norm{\nabla v}_3+\norm{\Lambda v}_3\le C\norm{v}_{\dot H^{3/2}},
\end{equation}
with mean-zero fields on the torus.
\end{lemma}
\begin{proof}
These are standard Sobolev estimates; for the Euclidean product bound
see Tao~\cite[Appendix~A, Lemma~A.8]{taodispersive}, together with the
$H^2$ embedding into $L^\infty$. We recall a short Fourier argument
that also covers the torus.
The elementary weight inequality
$\langle\xi\rangle^m\le C_m(\langle\eta\rangle^m+
\langle\xi-\eta\rangle^m)$ and the Fourier convolution formula give
\[
 \norm{vw}_{H^m}\le C_m\bigl(\norm{\widehat v}_1\norm{w}_{H^m}
                            +\norm{\widehat w}_1\norm{v}_{H^m}\bigr).
\]
Here the Fourier norm is $L^1$ in the whole space and $\ell^1$ on the
torus. The convolution bound follows by the triangle inequality for a
sum or integral of translates in $L^2$. Cauchy--Schwarz gives
$\norm{\widehat v}_1\le C\norm{v}_{H^2}$, since
$\int_{\R^3}\langle\xi\rangle^{-4}\dd\xi<\infty$ and
$\sum_{k\in\mathbb Z^3}\langle2\pi k\rangle^{-4}<\infty$.
The latter follows by grouping lattice points in dyadic shells; the
shell of radius $2^j$ contributes at most $C2^{-j}$.
This proves the product bounds, first for Schwartz functions or Fourier
polynomials and then by approximation. Absolute Fourier convergence
also proves the $L^\infty$ estimate. The argument applies componentwise
to vectors and tensors, and factorization gives the corresponding
product difference bound.

Lemma~\ref{lem:critical-embeddings} proves~\eqref{eq:embeddings} on both
domains. Its order-one case gives $H^1\hookrightarrow L^6$ on $\R^3$.
On the torus apply it to $v-\widehat v(0)$ and use
$|\widehat v(0)|\le\norm{v}_2$.
\end{proof}

\begin{proof}[Proof of Proposition~\ref{prop:local}]
\emph{Application of the classical local theory.}
Tao~\cite[Theorem~5.1(ii)--(iv)]{tao2013} gives local existence,
uniqueness, and smoothness for the forced periodic equation;
\cite[Theorem~5.4(ii)--(iv)]{tao2013} gives the whole-space counterpart.
On each finite interval our force is bounded into $H^1$ and integrable
into every $H^m$, so it satisfies the source hypotheses. The higher-order
bounds in part~(ii) hold on the same local interval for every order.
Moreover, the proof of Theorem~5.4(iv) explicitly allows
$a\in H^k$ and $f\in C^j_tH^k_x$ for all $j,k$, rather than requiring
Schwartz data. This gives the common-interval smoothness used here.
More explicitly, the bounds in every spatial Sobolev order and the
projected equation first give $u\in W^{1,\infty}_tH^k_x$ for each $k$,
hence continuity into every $H^k$. The equation then has a continuous
right-hand side in every $H^k$; repeated time differentiation gives
$C^j_tH^k_x$ regularity for all $j,k$, including one-sided derivatives
at the initial time. We apply the projected formulation with force $\PP f$; the pressure
formulas in Section~\ref{sec:analytic} restore the gradient part of $f$.

The cited theory uses viscosity one. For general $\nu>0$, set
\[
 \tau=\nu t,\qquad
 \widetilde u(\tau,x)=\nu^{-1}u(\tau/\nu,x),\qquad
 \widetilde p(\tau,x)=\nu^{-2}p(\tau/\nu,x),\qquad
 \widetilde f(\tau,x)=\nu^{-2}f(\tau/\nu,x).
\]
These fields satisfy the equation with viscosity one and preserve the
stated smoothness and Sobolev classes. Undoing this change restores $\nu$.

The periodic theorem is stated for mean-zero data. To remove this
restriction, define the known functions
\[
 m(t)=\int_{\TT}a\dd x+\int_0^t\int_{\TT}f(r,x)\dd x\dd r,
 \qquad X(t)=\int_0^t m(r)\dd r.
\]
The change of variables
\[
 v(t,x)=u(t,x+X(t))-m(t),\qquad
 h(t,x)=f(t,x+X(t))-m'(t)
\]
reduces the equation to mean-zero initial velocity and force, with
pressure $p(t,x+X(t))$. Indeed, the extra transport from $X'=m$
cancels the term subtracted from the advecting velocity. Spatial
translations preserve every Sobolev norm. Solving the reduced problem
and reversing the transformation therefore gives the periodic solution
with the required pressure normalization; compare
\cite[Section~3, equations~(36)--(38)]{tao2013}. No mean reduction is
needed on $\R^3$.

In either domain, the resulting velocity satisfies the mild equation
\begin{align}
 u(t)={}&e^{\nu t\Delta}a
 -\int_0^t e^{\nu(t-r)\Delta}\PP\nabla\cdot(u\otimes u)(r)\dd r\nonumber\\
 &+\int_0^t e^{\nu(t-r)\Delta}\PP f(r)\dd r.\label{eq:mild}
\end{align}
Conversely, smooth solutions in our classes satisfy this equation,
using the prescribed pressure gradient and the heat semigroup. For
completeness, if $z=u_1-u_2$ for two such solutions with the same data,
\[
 \tfrac12(\norm{z}_2^2)'+\nu\norm{\nabla z}_2^2
 \le\norm{\nabla u_2}_\infty\norm{z}_2^2.
\]
The coefficient is bounded on compact common intervals, since
$u_2\in C_tH^3$. Gr\"onwall gives uniqueness. Patching these local
solutions defines the maximal lifespan, as in
\cite[Corollary~5.2]{tao2013} and the corresponding whole-space argument.

\emph{The stated continuation criterion.}
We retain the short energy argument that yields precisely
\eqref{eq:criterion}. For every integer $m\ge3$, pairing the equation
with $u$ in $H^m$, integrating by parts, and using
\eqref{eq:Rproduct} gives
\begin{equation}\label{eq:Rhigh}
 \frac12\frac{\dd}{\dd t}\norm{u}_{H^m}^2
 +\nu\norm{\nabla u}_{H^m}^2
 \le C_m\norm{u}_{H^2}\norm{u}_{H^m}\norm{\nabla u}_{H^m}
      +\norm{f}_{H^m}\norm{u}_{H^m}.
\end{equation}
The pressure term vanishes by solenoidality. These identities are
justified by Fourier approximation on compact intervals of smooth
existence. Young's inequality and division by the regularized norm
$(\norm{u}_{H^m}^2+\zeta^2)^{1/2}$, followed by $\zeta\downarrow0$, imply
\begin{equation}\label{eq:highcontinuation}
 (\norm{u}_{H^m})'
 \le C_{m,\nu}\norm{u}_{H^2}^2\norm{u}_{H^m}+\norm{f}_{H^m}.
\end{equation}
If \eqref{eq:criterion} holds, Gr\"onwall bounds every $H^m$ norm
uniformly up to $S$. The $H^1$ local existence bounds of the cited
Theorems~5.1(ii) and~5.4(ii) then give a common positive existence
duration when restarting at $t_0\uparrow S$: the initial $H^1$ norms
stay bounded, and $f$ is bounded into $H^1$ on $[0,S+1]$.
Their higher-order regularity statements give smooth solutions on
these same intervals. One interval extends beyond $S$, and uniqueness
identifies it with the original solution on the overlap. The periodic
mean stays bounded by its displayed formula, so the same restart
argument applies there. This proves the asserted continuation without
requiring zero mean or a whole-space Poincar\'e inequality.
\end{proof}
